% !TeX program = xelatex
% 从本文件所在目录编译；复制模板时一并复制 sty、book-path.tex 和 assets。
\documentclass[aspectratio=169,10pt]{beamer}
\usepackage[UTF8,fontset=fandol]{ctex}
% 优先匹配原模板字体，其他系统自动使用 TeX Live 自带的 Fandol。
\IfFontExistsTF{Microsoft YaHei}{\setCJKsansfont{Microsoft YaHei}}{}
\usetheme{GaoZhongYuan}

\title{标题}
\author{姓名}
\date{2026.08.12}

\begin{document}
\campusmaketitle

\begin{frame}[t]{\underline{任务一：}}
  % 在这里填写正文，支持公式、列表、图片和 Beamer 分步显示。
  % 例如：\[ (a+b)^2=a^2+2ab+b^2. \]
\end{frame}

% 练习题一题一页；每个小问单独复制一个环境。
\begin{campuspractice}{课堂练习}
  已知 $\log_2x=3$，求 $x$。
\end{campuspractice}
% 解答另起一页，不占用演算区。
\begin{frame}[t]{课堂练习：参考解答}
  $\log_2x=3\Longleftrightarrow x=2^3=8$。
\end{frame}
\begin{campussummary}
  \begin{itemize}
    \item 本课核心概念与适用条件。
    \item 典型方法及其推理依据。
    \item 易错点与检查方法。
  \end{itemize}
\end{campussummary}
% 复制上面的环境添加题目；结束前保留课堂小结。
\campusclosing
\end{document}
